\section{Augmented Hierarchical Model, the Rigid-Link Predictive Density, and Mean-Field Fixed Points}
\label{app:augmented_joint}

The main text reads the cross-scale shadow $p_i^{(\ell)} = \Omega_{i,I}[q_I^{(\ell+1)}]$ of Section~\ref{sec:cross_scale_shadows} as the one-episode-lagged predictive prior of an episode-indexed joint. When a parent forms at episode $n$, the prior of its children at episode $n+1$ is the filtered parent law pushed through a declared transition, and for a transported-copy transition this is the shadow formula exactly~\cite[Ch.~12]{DennisGaugeVFE2026}. This appendix records what a same-time hierarchical joint does and does not supply. We define an augmented hierarchical generative model with Gaussian link kernels on the tree of agent latents, prove that it is a proper joint, state its mean-field fixed-point equations, and compute the variational predictive density whose rigid-link limit reproduces the shadow formula. That density is not a message of the joint: evaluated against the full parent posterior it reuses each constituent's own evidence, and the functional built from it is not an evidence bound (Remark~\ref{rem:shadow_marginal_vs_mf}). The general statements behind this are the typing prohibition of the companion theory, under which no generative kernel reads a recognition law, and its result that an energy appended to a joint without the corresponding normalizer is not an evidence bound~\cite[Ch.~4]{DennisGaugeVFE2026}. The companion theory's downward kernels and conditional cross-scale divergence give the evidence-consistent form of top-down influence~\cite[Ch.~12]{DennisGaugeVFE2026}.

\subsection{The Augmented Joint}

Let the hierarchy be a finite tree of depth $\ell_{\max}$ and finite breadth, with $I_\ell$ the set of agents at level $\ell$ and $\pi(i)$ the parent of agent $i$. Each agent $i \in I_\ell$ carries a latent variable $k_i^{(\ell)} \in \mathbb{R}^{K_q}$. The construction is made at a single base point, so it is a finite-dimensional probability law on the product space $\mathbb{K} := \prod_{\ell, i \in I_\ell} \mathbb{R}^{K_q}$.

\begin{definition}[Augmented hierarchical joint]
\label{def:augmented_joint}
For each parent-child edge $(i, \pi(i))$ with $i \in I_\ell$, define the Gaussian link kernel
\begin{equation}
\mathcal{K}_{i,\pi(i)}^{(\ell)}\big(k_i^{(\ell)} \mid k_{\pi(i)}^{(\ell+1)}\big)
=
(2\pi\sigma^2)^{-K_q/2} \exp\left[-\frac{1}{2\sigma^2}
\big\|k_i^{(\ell)} - \Omega_{i,\pi(i)} k_{\pi(i)}^{(\ell+1)}\big\|^2\right],
\label{eq:gibbs_factor}
\end{equation}
where $\Omega_{i,\pi(i)} = \exp(\phi_i^{(\ell)})\exp(-\phi_{\pi(i)}^{(\ell+1)}) \in G$ is the inter-level transport of Section~\ref{sec:transport_hierarchy} and $\sigma^2 > 0$ is a link-stiffness parameter. The kernel is a normalized transition density, a propagator from parent to child, and not a Boltzmann weight. The top-level prior $p_{\mathrm{top}}(k^{(\ell_{\max})})$ is a declared root, fixed as a boundary condition (Section~\ref{sec:ouroboros}). The augmented joint over all agent latents is
\begin{equation}
p_{\sigma}\big(\{k_i^{(\ell)}\}, o\big)
=
p(o \mid \{k_i^{(0)}\}) \prod_{\ell = 0}^{\ell_{\max}-1} \prod_{i \in I_\ell} \mathcal{K}_{i,\pi(i)}^{(\ell)}\big(k_i^{(\ell)} \mid k_{\pi(i)}^{(\ell+1)}\big) p_{\mathrm{top}}\big(k^{(\ell_{\max})}\big).
\label{eq:augmented_joint}
\end{equation}
\end{definition}

\begin{lemma}[Well-definedness]
\label{lem:aug_joint_welldefined}
Assume the observation likelihood $p(o \mid \cdot)$ is bounded above by a finite constant (unity in the discrete-observation case, the normalization constant in the Gaussian case). Then for every finite $\sigma^2 > 0$ and every finite hierarchical tree, the augmented joint~\eqref{eq:augmented_joint} is integrable over $\mathbb{K}$ for every fixed $o$, with normalization equal to the evidence $p(o)$; the prior portion, obtained by omitting the likelihood factor, is a proper probability density on $\mathbb{K}$.
\end{lemma}

\begin{proof}
Each kernel $\mathcal{K}_{i,\pi(i)}^{(\ell)}$ is a Gaussian density on $k_i^{(\ell)}$ with mean $\Omega_{i,\pi(i)} k_{\pi(i)}^{(\ell+1)}$ and covariance $\sigma^2 I$, and it integrates to unity for every value of $k_{\pi(i)}^{(\ell+1)}$. The factor graph of~\eqref{eq:augmented_joint} is the hierarchical tree itself. For tree-structured products of proper conditional densities, the integral over the leaves factorizes through the tree~\cite[\S3.2]{Lauritzen1996} and leaves the top-level prior $p_{\mathrm{top}}$, which is normalized, so the prior portion is a proper probability density on $\mathbb{K}$. The likelihood enters as a multiplicative factor at the leaves and is bounded above by hypothesis, so it does not affect integrability, and integrating it against the prior portion yields $p(o)$.
\end{proof}

\subsection{Mean-Field Fixed-Point Equations}

Mean-field variational inference on the augmented joint factorizes the variational posterior across agent latents.

\begin{lemma}[Mean-field fixed-point equations]
\label{lem:mf_fixedpoint}
The variational free energy
\begin{equation}
\mathcal{F}_{\sigma}[q]
=
\sum_{\ell, i} \mathbb{E}_{q_i^{(\ell)}}\big[\log q_i^{(\ell)}\big]
- \sum_{\ell < \ell_{\max}, i} \mathbb{E}_{q_i^{(\ell)} q_{\pi(i)}^{(\ell+1)}}\big[\log \mathcal{K}_{i,\pi(i)}^{(\ell)}\big]
- \mathbb{E}_{q^{(\ell_{\max})}}\big[\log p_{\mathrm{top}}\big]
- \mathbb{E}_{q^{(0)}}\big[\log p(o \mid k^{(0)})\big]
\label{eq:mf_free_energy}
\end{equation}
is stationary under the mean-field factorization $q(\{k_i^{(\ell)}\}) = \prod_{\ell, i} q_i^{(\ell)}(k_i^{(\ell)})$ when each factor satisfies
\begin{equation}
\begin{aligned}
\log q_i^{(\ell)}(k_i)
={}&
\text{const}
+ \mathbb{E}_{q_{\pi(i)}^{(\ell+1)}}\big[\log \mathcal{K}_{i,\pi(i)}^{(\ell)}(k_i \mid k_{\pi(i)})\big]
+ \sum_{j : \pi(j) = i} \mathbb{E}_{q_j^{(\ell-1)}}\big[\log \mathcal{K}_{j,i}^{(\ell-1)}(k_j \mid k_i)\big]\\
&+ \delta_{\ell,0} \mathbb{E}_{q^{(0)}_{\setminus i}}\big[\log p(o \mid k^{(0)})\big]
+ \delta_{\ell,\ell_{\max}} \log p_{\mathrm{top}}(k_i),
\end{aligned}
\label{eq:mf_fixedpoint}
\end{equation}
where $q^{(0)}_{\setminus i}$ is the product of the other leaf factors. The first expectation is the top-down message from the parent of agent $i$ (absent at $\ell = \ell_{\max}$), the sum collects the bottom-up messages from its children (absent at $\ell = 0$), and the data likelihood enters only at the leaves.
\end{lemma}

\begin{proof}
Equation~\eqref{eq:mf_fixedpoint} is the standard coordinate-ascent mean-field update~\cite[\S5.3]{Wainwright2008}, \cite[\S10.1]{Bishop2006}. It follows by varying~\eqref{eq:mf_free_energy} with respect to $q_i^{(\ell)}$ at fixed values of the other factors, imposing $\int q_i^{(\ell)} = 1$ by a Lagrange multiplier, and collecting the terms that depend on $k_i^{(\ell)}$.
\end{proof}

\subsection{The Rigid-Link Variational Predictive Density}

\begin{lemma}[Rigid-link variational predictive density]
\label{lem:shadow_mf_optimum}
Let the variational posterior of the parent be Gaussian, $q_{\pi(i)}^{(\ell+1)} = \mathcal{N}(\mu_{\pi(i)}^{(\ell+1)}, \Sigma_{\pi(i)}^{(\ell+1)})$. The variational predictive density of agent $i$, the marginal of the link kernel against the parent's variational posterior, is
\begin{equation}
\hat{p}_i^{(\ell),\sigma}(k_i) := \int \mathcal{K}_{i,\pi(i)}^{(\ell)}(k_i \mid k_{\pi(i)}) q_{\pi(i)}^{(\ell+1)}(k_{\pi(i)}) dk_{\pi(i)}
= \mathcal{N}\big(k_i; \Omega_{i,\pi(i)} \mu_{\pi(i)}^{(\ell+1)}, \Omega_{i,\pi(i)} \Sigma_{\pi(i)}^{(\ell+1)} \Omega_{i,\pi(i)}^\top + \sigma^2 I\big),
\label{eq:shadow_soft}
\end{equation}
and as $\sigma^2 \to 0$ it converges to the pushforward $(\Omega_{i,\pi(i)})_* q_{\pi(i)}^{(\ell+1)}$, which is the shadow formula of Eq.~\eqref{eq:cross_scale_shadow}. The lemma computes a density. That density is neither the mean-field factor of Eq.~\eqref{eq:mf_fixedpoint} nor a message of the joint~\eqref{eq:augmented_joint} (Remark~\ref{rem:shadow_marginal_vs_mf}).
\end{lemma}

\begin{proof}
Conditional on $k_{\pi(i)}$, the kernel~\eqref{eq:gibbs_factor} is Gaussian with mean $\Omega_{i,\pi(i)} k_{\pi(i)}$, linear in the parent, and covariance $\sigma^2 I$. The linear-Gaussian marginalization rule~\cite[\S2.3.3]{Bishop2006} gives a Gaussian marginal with mean $\Omega_{i,\pi(i)}\mu_{\pi(i)}^{(\ell+1)}$ and covariance $\Omega_{i,\pi(i)}\Sigma_{\pi(i)}^{(\ell+1)}\Omega_{i,\pi(i)}^\top + \sigma^2 I$. As $\sigma^2 \to 0$ the kernel converges weakly to $\delta(k_i - \Omega_{i,\pi(i)} k_{\pi(i)})$, and the marginal converges to
\begin{equation}
\int \delta\big(k_i - \Omega_{i,\pi(i)} k_{\pi(i)}\big) q_{\pi(i)}^{(\ell+1)}(k_{\pi(i)}) dk_{\pi(i)} = \big(\Omega_{i,\pi(i)}\big)_* q_{\pi(i)}^{(\ell+1)}.
\end{equation}
\end{proof}

\begin{remark}[The predictive density is not a message of the joint]
\label{rem:shadow_marginal_vs_mf}
Four facts fix the status of Lemma~\ref{lem:shadow_mf_optimum}. First, the strict mean-field factor is a different density. Substituting the kernel into the top-down term of Eq.~\eqref{eq:mf_fixedpoint} gives
\[
\mathbb{E}_{q_{\pi(i)}^{(\ell+1)}}\big[\log \mathcal{K}_{i,\pi(i)}^{(\ell)}\big] = \text{const} - \frac{1}{2\sigma^2}\Big[\big\|k_i - \Omega_{i,\pi(i)}\mu_{\pi(i)}^{(\ell+1)}\big\|^2 + \operatorname{tr}\big(\Omega_{i,\pi(i)}\Sigma_{\pi(i)}^{(\ell+1)}\Omega_{i,\pi(i)}^\top\big)\Big],
\]
so in isolation from children and data the mean-field factor is $\mathcal{N}(\Omega_{i,\pi(i)}\mu_{\pi(i)}^{(\ell+1)}, \sigma^2 I)$: the parent covariance enters only the constant, and in the rigid limit the factor collapses to a point mass. Second, the predictive density~\eqref{eq:shadow_soft} equals the prior marginal of $k_i$ under $p_\sigma$ when $q_{\pi(i)}^{(\ell+1)}$ is the prior marginal of the parent. A parent posterior, by contrast, has absorbed the evidence of every leaf below the parent, including the leaves below agent $i$.

Third, evaluated against the full parent posterior, the shadow reuses each constituent's own evidence; the evidence-consistent downward message uses the parent posterior with that constituent's contribution removed. On a tree, belief propagation is exact~\cite[\S8.4]{Bishop2006}: the exact posterior of $k_i$ is proportional to the bottom-up evidence of agent $i$ times the cavity predictive density $\int \mathcal{K}_{i,\pi(i)}^{(\ell)}(k_i \mid k_{\pi(i)}) q_{\pi(i)}^{\setminus i}(k_{\pi(i)}) dk_{\pi(i)}$, where $q_{\pi(i)}^{\setminus i}$ is the parent posterior computed without the message from agent $i$. Multiplying the full-posterior predictive density by agent $i$'s likelihood counts that likelihood twice. In a Gaussian instance with one parent, two children, $K_q = 2$, $\mathrm{GL}^+(2)$ transports, and Gaussian observations under each child, the cavity construction reproduces the exact child posterior to $10^{-16}$ at $\sigma^2 = 0.3$, while the full-posterior shadow scheme errs in the child mean by $0.13$; at $\sigma^2 = 10^{-8}$ its precision exceeds the exact child precision by exactly the precision of the child's own observation.

Fourth, the functional built from the predictive density is not an evidence bound. Replacing the mean-field cross term $\mathbb{E}_{q_i q_{\pi(i)}}[-\log \mathcal{K}_{i,\pi(i)}]$ by the predictive cross-entropy $\mathbb{E}_{q_i}[-\log \hat{p}_i]$ can only lower the objective, since $-\log$ is convex and Jensen's inequality gives $\mathbb{E}_{q_i}[-\log \int \mathcal{K}_{i,\pi(i)} q_{\pi(i)}] \le \mathbb{E}_{q_i q_{\pi(i)}}[-\log \mathcal{K}_{i,\pi(i)}]$. The resulting same-time shadow functional $D_{\mathrm{KL}}(q_{\mathrm{top}}\|p_{\mathrm{top}}) + \sum_i [D_{\mathrm{KL}}(q_i\|\hat{p}_i) - \mathbb{E}_{q_i}\log p(o_i \mid k_i)]$ can therefore fall below the surprisal $-\log p_\sigma(o)$, which no variational free energy of $p_\sigma$ can do. In a scalar instance with $p_{\mathrm{top}} = \mathcal{N}(0,1)$, $\sigma^2 = 0.05$, and observation variance $0.2$, its minimum over Gaussian factors lies $0.14$ nats below the surprisal for the concordant data $o = (1,1)$ and $12.98$ nats below for $o = (2,-2)$. In the rigid limit the child-parent term becomes $D_{\mathrm{KL}}(q_i\|(\Omega_{i,\pi(i)})_* q_{\pi(i)})$ with both factors varying. Its mixed third variation~\eqref{eq:consensus_third_variation} is nonzero, so Theorem~\ref{thm:state_level_elbo_nogo} applies to it exactly as to the peer terms, and the same-time shadow edge has no fixed state-level evidence-bound representation.
\end{remark}

\subsection{Evidence-Consistent Readings of Top-Down Influence}

Lemmas~\ref{lem:aug_joint_welldefined}--\ref{lem:shadow_mf_optimum} do not make the same-time shadow a derived prior. Four readings of top-down influence are consistent with one normalized joint. The first is the conditional coupling $\mathbb{E}_{q_{\pi(i)}} D_{\mathrm{KL}}(q_i\|\mathcal{K}_{i,\pi(i)}(\cdot \mid k_{\pi(i)}))$, the mean-field cross term of Eq.~\eqref{eq:mf_free_energy}; it is the conditional cross-scale divergence of the companion theory's downward-coupling proposition~\cite[Ch.~12]{DennisGaugeVFE2026}, and its rigid limit diverges as $1/\sigma^2$ unless the factors degenerate. The second is the cavity predictive density of Remark~\ref{rem:shadow_marginal_vs_mf}, the exact downward message on a tree. The third is the random-effects, or empirical-Bayes, reading: the parent is a point hyperparameter $\theta = (m, S)$, the constituents are conditionally independent draws $k_i \sim \mathcal{N}(\Omega_{i,I} m, \Omega_{i,I} S \Omega_{i,I}^\top)$, and the shadow is then an honest prior. Fitting $\theta$ is type-II maximum likelihood, the barycenter of Eq.~\eqref{eq:meta_agent_barycenter} with its dispersion term is the exact update of $\theta$ at fixed constituent factors~\cite{Neal1998}, and the parent carries no belief of its own; $\Sigma_I$ is then a population covariance and not a posterior uncertainty, so this is a different model from Definition~\ref{def:augmented_joint}. The fourth is the one-episode-lagged reading of the main text, stated at the opening of this appendix, with the transported-copy transition $k_i^{(n+1)} = \Omega_{i,I} k_I^{(n)}$; Eq.~\eqref{eq:cross_scale_shadow} is its formula with the variational parent standing in for the filtered law. The first three readings are same-time; the fourth carries the participatory reading of Section~\ref{sec:cross_scale_shadows}.

In hierarchical generative models of cortical inference~\cite{friston2008hierarchical, parr2022active} the level-$\ell$ prior is the conditional $p(x^{(\ell)} \mid x^{(\ell+1)})$ of the generative model; the link kernel~\eqref{eq:gibbs_factor} plays that role here, with a gauge-equivariant mean and an isotropic variance. It is also a special case of the top-down Gaussian conditional of hierarchical variational autoencoders~\cite{Sonderby2016}, whose conditional variances are learned per dimension. Their bound evaluates this conditional at parent samples, which is the conditional coupling above and not the predictive density of Lemma~\ref{lem:shadow_mf_optimum}, and the rigid limit $\sigma^2 \to 0$ is singular for that variational problem.

The shadow formula of the main text is the rigid-link limit $\sigma^2 = 0$. The finite-$\sigma^2$ density~\eqref{eq:shadow_soft} inflates the transported parent covariance by $\sigma^2 I$ and can regularize an ill-conditioned $\Omega_{i,I}\Sigma_I\Omega_{i,I}^\top$. It defines a one-parameter family of softer shadows for future study; no simulation at small positive $\sigma^2$ is reported.

The lemmas concern the cross-scale terms alone. Reinstating the within-scale alignment terms gives a fixed-point system in which each level-$\ell$ posterior balances the transported parental pull against the within-scale alignment energy. Because the softmax weights $\beta_{ij}$ are row-stochastic, a ``small $\beta_{ij}$'' is not a small parameter. The dimensionless coupling strengths are the raw alignment weights $\tilde\beta_{ij}, \tilde\gamma_{ij}$ of Section~\ref{sec:raw_alignment_weights}, and the within-scale terms vanish from the functional in the limit $\tilde\beta_{ij}, \tilde\gamma_{ij} \to 0$. Perturbative control of the within-scale corrections is open (Section~\ref{app:mf_fixedpoint_status}).

\subsection{Status of the Mean-Field Fixed-Point Question}
\label{app:mf_fixedpoint_status}

At zero within-scale coupling, mean field on the augmented joint is exact and uniquely solvable (Section~\ref{app:cross_scale_gaussian_exception}). The shadow dynamics of Section~\ref{sec:cross_scale_shadows} are not this mean field (Remark~\ref{rem:shadow_marginal_vs_mf}) and inherit none of these guarantees. The within-scale question has three separate answers: a general obstruction for a fixed joint on the original state variables, a self-consistent auxiliary model that becomes fixed only after an equilibrium is selected, and an exact Gibbs lift whose random variable is an entire belief configuration (Theorem~\ref{thm:belief_configuration_gibbs_lift}, Appendix~\ref{sec:meta_agent_rg}).

\subsubsection{Fixed Representation on the State Variables: General Obstruction}

Let $(\mathsf{X}_i,\nu_i)$ be the state space of agent $i$, and let $\mathsf{T}_{ij}$ denote the fixed mass-preserving positive linear transport induced by $(\Omega_{ij})_*$. A family $\mathcal{Q}_i$ of strictly positive densities is called mixture-open here when every zero-mass tangent $h_i$ under consideration with $h_i/q_i\in L^\infty(q_i)$ generates $q_i(t)=q_i+t h_i\in\mathcal{Q}_i$ for all sufficiently small $|t|$. All differentiated integrals below are assumed finite.

\begin{theorem}[State-level mean-field ELBO obstruction]
\label{thm:state_level_elbo_nogo}
Let the factors $q_i\in\mathcal{Q}_i$ vary independently, let at least one nonself coefficient $\beta_{ij}>0$, and consider
\begin{equation}
\mathcal{C}_{\mathrm{cons}}(q)
=\mathcal{F}_{\mathrm{rest}}(q)
+\sum_{a\ne b}\beta_{ab}D_{\mathrm{KL}}(q_a\|\mathsf{T}_{ab}q_b),
\label{eq:population_consensus_nogo_functional}
\end{equation}
where $D^3_{q_iq_jq_j}\mathcal{F}_{\mathrm{rest}}=0$ for $i\ne j$. Assume that for one active edge $(i,j)$ there is an admissible tangent $h_j$ such that
\begin{equation}
g_{ij}:=\left(\frac{\mathsf{T}_{ij}h_j}{\mathsf{T}_{ij}q_j}\right)^2
\label{eq:nogo_nondegenerate_tangent}
\end{equation}
is bounded and nonconstant $q_i$-almost everywhere. Then there is no density $p_\theta$ on $\prod_i\mathsf{X}_i$, independent of $q$, and no constant $C_0$ such that
\begin{equation}
\mathcal{C}_{\mathrm{cons}}(q)
=D_{\mathrm{KL}}\left(\bigotimes_iq_i\middle\|p_\theta\right)+C_0,
\label{eq:state_level_elbo_representation_impossible}
\end{equation}
throughout the open product family.
\end{theorem}

\begin{proof}
Write $Q_q=\bigotimes_iq_i$. The standard fixed-model mean-field decomposition~\cite{Wainwright2008} gives
\begin{equation}
D_{\mathrm{KL}}(Q_q\|p_\theta)
=\sum_i\int_{\mathsf{X}_i}q_i\log q_i d\nu_i
-\int_{\prod_i\mathsf{X}_i}\left(\prod_iq_i\right)\log p_\theta d\nu_{1:N}.
\label{eq:fixed_joint_mean_field_structure}
\end{equation}
The entropy is separable, while the fixed-model energy is affine in each factor. Therefore
\begin{equation}
D^3_{q_iq_jq_j}D_{\mathrm{KL}}(Q_q\|p_\theta)[a_i,h_j,h_j]=0,
\qquad i\ne j.
\label{eq:fixed_joint_third_variation_zero}
\end{equation}
For the active consensus edge, take $q_i(s)=q_i+s a_i$ and $q_j(t)=q_j+t h_j$. Direct differentiation gives
\begin{equation}
\left.
\frac{\partial^3}{\partial s\partial t^2}
D_{\mathrm{KL}}(q_i(s)\|\mathsf{T}_{ij}q_j(t))
\right|_{s=t=0}
=\int_{\mathsf{X}_i}a_i g_{ij}d\nu_i.
\label{eq:consensus_third_variation}
\end{equation}
Choose $a_i=q_i(g_{ij}-\mathbb{E}_{q_i}[g_{ij}])$, which has zero mass and satisfies $a_i/q_i\in L^\infty(q_i)$ because $g_{ij}$ is bounded; it is therefore admissible for a sufficiently small mixture path. Equation~\eqref{eq:consensus_third_variation} becomes $\operatorname{Var}_{q_i}(g_{ij})>0$, contradicting Eq.~\eqref{eq:fixed_joint_third_variation_zero}. The terms collected in $\mathcal{F}_{\mathrm{rest}}$ cannot cancel this mixed variation by hypothesis. Every other ordered-edge term also has zero $q_i,q_j,q_j$ variation: an unrelated edge omits one varied factor, while the reverse edge is affine in $q_j$ outside its entropy and that entropy's second variation is independent of $q_i$.
\end{proof}

\begin{remark}[Scope of the obstruction]
\label{rem:state_level_elbo_nogo_scope}
The theorem applies directly to the entropy-retaining functional by holding its independently varied attention rows fixed. The self terms, likelihood terms with fixed observation models, and categorical attention entropies lie in $\mathcal{F}_{\mathrm{rest}}$. The same argument applies to the child-parent edge of a same-time shadow (Remark~\ref{rem:shadow_marginal_vs_mf}). The result concerns an ordinary mean-field factorization on the original agent-state variables. It does not exclude frozen source templates, restricted quadratic families, compatible auxiliary variables, a model selected after a fixed point, or a probability law over belief configurations. If the attention optimum $\beta^*(q)$ has already been substituted, its response derivatives define a different reduced functional and require a separate representation test.
\end{remark}

\subsubsection{Exact Cross-Scale Gaussian Exception}
\label{app:cross_scale_gaussian_exception}

At zero raw within-scale coupling ($\tilde\beta_{ij}=\tilde\gamma_{ij}=0$, so the alignment sectors are absent), the augmented joint of Definition~\ref{def:augmented_joint} is a closed variational object. Assume the top-level prior $p_{\mathrm{top}}$ is a nondegenerate Gaussian with full-rank precision $\Lambda_{\mathrm{top}}$, and take the leaf observation term absent or Gaussian; this is needed both for the negative log joint to be a quadratic form in the stacked latents and for the positive-definiteness below. The precision matrix of that quadratic form is block-sparse along the hierarchical tree, with nonzero blocks only on the diagonal and on parent-child edges. The diagonal block of agent $i$ is $\sigma^{-2}\big(I + \sum_{j : \pi(j) = i}\Omega_{j,i}^\top\Omega_{j,i}\big)$, where the identity contribution is absent at the top level (the prior precision $\Lambda_{\mathrm{top}}$ takes its place) and the leaf blocks additionally carry any Gaussian likelihood curvature; the off-diagonal block on each parent-child edge is $-\sigma^{-2}\Omega_{i,\pi(i)}$ together with its transpose. The link terms alone are only positive semidefinite, with kernel the consistent-propagation subspace on which $k_i = \Omega_{i,\pi(i)}k_{\pi(i)}$ along every edge. The full-rank prior precision at the top removes this kernel, so for finite $\sigma^2 > 0$ the stacked precision is symmetric positive definite, the joint is a proper Gaussian, and the coordinate-ascent mean-field updates of Lemma~\ref{lem:mf_fixedpoint} are Gauss-Seidel sweeps on a symmetric positive-definite system, which converge to the unique global optimum. In this regime the evidence-bound interpretation is exact and the fixed point exists and is unique. For a single parent and its constituents, an anchored Gaussian star, the companion supplement gives the exact fixed point and the contraction rate of this coordinate-ascent schedule, with Gaussian observations of the constituents included~\cite[Ch.~G10]{DennisGaugeVFESuppG2026}.

\subsubsection{Self-Consistent Equilibrium Representation}

\begin{theorem}[Existence of a self-consistent source model]
\label{thm:self_consistent_source_equilibrium}
Let the latent alphabet $\mathsf{X}$ and the agent set be finite. Let $p_i(k)>0$, let $\pi_{ij}>0$ with $\sum_j\pi_{ij}=1$, and let $a_i>0$, $\lambda>0$, and $\tau>0$. Assume every stochastic transport has a uniform floor: there is $\epsilon>0$ such that $(\mathsf{T}_{ij}q)(k)\ge\epsilon$ for every probability vector $q$, state $k$, and active edge. Define
\begin{equation}
E_{ij}(q):=D_{\mathrm{KL}}(q_i\|\mathsf{T}_{ij}q_j),
\qquad
\beta_{ij}(q):=
\frac{\pi_{ij}\exp[-E_{ij}(q)/\tau]}
{\sum_m\pi_{im}\exp[-E_{im}(q)/\tau]},
\label{eq:self_consistent_attention}
\end{equation}
and $\bar{a}_i:=a_i+\lambda$. The geometric-pooling response map
\begin{equation}
B_i(q)(k):=
\frac{
p_i(k)^{a_i/\bar{a}_i}
\prod_j[\mathsf{T}_{ij}q_j(k)]^{\lambda\beta_{ij}(q)/\bar{a}_i}
}{
\sum_{u\in\mathsf{X}}
p_i(u)^{a_i/\bar{a}_i}
\prod_j[\mathsf{T}_{ij}q_j(u)]^{\lambda\beta_{ij}(q)/\bar{a}_i}
}
\label{eq:self_consistent_response_map}
\end{equation}
has at least one receiver-response equilibrium $q^\star=B(q^\star)$. Write $\beta^\star:=\beta(q^\star)$.

At any such fixed point, define normalized row distributions
\begin{equation}
P_i^\star(k,j):=\pi_{ij}[\mathsf{T}_{ij}q_j^\star](k),
\qquad
Q_i^\star(k,j):=\beta_{ij}^\star q_i^\star(k).
\label{eq:equilibrium_source_joint}
\end{equation}
They satisfy the exact identity
\begin{equation}
D_{\mathrm{KL}}(Q_i^\star\|P_i^\star)
=\sum_j\beta_{ij}^\star D_{\mathrm{KL}}(q_i^\star\|\mathsf{T}_{ij}q_j^\star)
+D_{\mathrm{KL}}(\beta_i^\star\|\pi_i).
\label{eq:equilibrium_row_kl_identity}
\end{equation}
The identity is algebraic: it holds with $(q^\star,\beta^\star)$ replaced by any $(q,\beta)$ in which each $q_i$ and each row $\beta_i$ is a probability vector, and in particular at the equilibrium.
\end{theorem}

\begin{proof}
The uniform floor makes every energy in Eq.~\eqref{eq:self_consistent_attention} continuous on the closed product simplex. Hence $\beta(q)$ and $B(q)$ are continuous, and $B$ maps that compact convex simplex into itself. Brouwer's theorem gives a receiver-response equilibrium. For fixed source templates and fixed $\beta_i$, the Euler equation for
\begin{equation}
a_iD_{\mathrm{KL}}(r\|p_i)
+\lambda\sum_j\beta_{ij}D_{\mathrm{KL}}(r\|\mathsf{T}_{ij}q_j),
\label{eq:frozen_row_geometric_pool_objective}
\end{equation}
under $\sum_k r(k)=1$ gives $r=B_i(q)$, so the fixed point is self-consistent with the frozen-row optimum. Expanding $\log[Q_i/P_i]=\log(\beta_{ij}/\pi_{ij})+\log[q_i(k)/\mathsf{T}_{ij}q_j(k)]$ and summing against $Q_i$ over $(k,j)$, using $\sum_kq_i(k)=\sum_j\beta_{ij}=1$, proves Eq.~\eqref{eq:equilibrium_row_kl_identity} for every such $(q,\beta)$.
\end{proof}

\begin{remark}[What the equilibrium construction proves]
\label{rem:equilibrium_model_scope}
The distributions $P_i^\star$ are fixed and normalized only after $q^\star$ has been selected. Their product is a proper auxiliary model over independent row variables, but it duplicates the transported source templates. It is not automatically a joint over shared latent variables $(k_1,\ldots,k_N)$. Such a joint additionally requires compatible conditionals, an acyclic factorization, or an undirected normalized potential. Equation~\eqref{eq:equilibrium_row_kl_identity} matches the ordinary row KL at unit attention temperature; arbitrary $\tau$ gives the tempered loss-based update of Eq.~\eqref{eq:self_consistent_attention}, a generalized-Bayesian construction and not one unscaled KL~\cite{BissiriHolmesWalker2016}. Because the response holds incoming source roles fixed, its equilibrium is not thereby a stationary point of the global smoothing scalar. The theorem proves existence, not uniqueness or convergence of a particular update schedule.
\end{remark}

\subsubsection{Scalar Objectives and Update Fields}

The within-scale construction requires a distinction between three objects. First, the canonical entropy-retaining row objective is $\mathcal{F}_{i,\mathrm{full}}=\sum_j\beta_{ij}E_{ij}+\tau\sum_j\beta_{ij}\log(\beta_{ij}/\pi_{ij})$. At its simplex optimum it is exactly $-\tau\log Z_i$, and the envelope theorem gives $\mathrm{d}\mathcal{F}_{i,\mathrm{red}}=\sum_j\beta_{ij}^*\mathrm{d}E_{ij}$. Summing over rows therefore gives one scalar potential with a symmetric Hessian. Second, the entropy-suppressed quantity $S=\sum_{ij}\beta_{ij}^*(x)E_{ij}(x)$ is a different scalar. Its response terms change its gradient but do not change the symmetry of its Hessian. Third, the deliberately asymmetric field
\begin{equation}
V_i^{\mathrm{rec}}(x):=\partial_{x_i}\left(\sum_j\beta_{ij}^*(x)E_{ij}(x)\right),
\qquad
\text{with all }\partial_{x_i}\left(\sum_m\beta_{\ell m}^*E_{\ell m}\right),\ \ell\ne i,\ \text{discarded},
\label{eq:receiver_only_field}
\end{equation}
is not the gradient of either scalar and is generically nonintegrable because cross-row mixed derivatives need not match. This nonpotential field is an optional update approximation, not a property of the softmax, the noncompact group, or differentiation of a scalar. The consensus functional remains a composite scalar potential and not a fixed state-level generative-model evidence bound; that distinction does not require a nonsymmetric Hessian.

Existence for the full Gaussian system is more delicate than the compact finite-state result of Theorem~\ref{thm:self_consistent_source_equilibrium}. A restricted stationary point exists under a symmetry hypothesis: the mean-only consensus state $\mu_i \equiv 0$ is stationary for every prior strength $\alpha>0$, temperature $\tau>0$, and finite $N$ when the prior means vanish and observation forcing is absent, because every pairwise alignment mean-gradient $\Lambda_{t,ij}(\mu_i-\Omega_{ij}\mu_j)$ and the prior pull vanish there. This does not prove a fixed point for the covariance and frame coordinates. For live state-dependent softmax weights on $\mathrm{GL}^+(K_q)$, local stability, global uniqueness, and global convergence remain open, and noncompact frame directions prevent a direct compactness argument unless one imposes a gauge slice, a bounded invariant container, or a coercive frame prior. A Banach guarantee holds in a restricted setting. Freeze covariances and attention weights, restrict transports to the orthogonal group, and take isotropic couplings $\Lambda_{t,ij}=\lambda_{ij}I$ with an arbitrary symmetric positive-definite $\Lambda_p$. The mean sub-map $\mu_i\leftarrow(\alpha\Lambda_p+b_iI)^{-1}(\alpha\Lambda_p\mu_{p,i}+\sum_j\beta_{ij}\lambda_{ij}\Omega_{ij}\mu_j)$, with $b_i=\sum_j\beta_{ij}\lambda_{ij}$, then has row norm $\sum_j\|(\alpha\Lambda_p+b_iI)^{-1}\beta_{ij}\lambda_{ij}\Omega_{ij}\|_2=b_i/[\alpha\lambda_{\min}(\Lambda_p)+b_i]<1$, so it is a contraction in the norm $\max_i\|\mu_i\|_2$. The isotropy hypothesis cannot be dropped. For anisotropic $\Lambda_{t,ij}$ that do not commute with their weighted sum, the operator $(\alpha\Lambda_p+\sum_j\beta_{ij}\Lambda_{t,ij})^{-1}\beta_{ij}\Lambda_{t,ij}$ can have norm above $\beta_{ij}\lambda_{\max}(\Lambda_{t,ij})/\lambda_{\max}(\sum_j\beta_{ij}\Lambda_{t,ij})$. In an instance with $N=3$, $K_q=2$, $\alpha=0.05$, $\Lambda_p=I$, anisotropic $\Lambda_{t,ij}$, and rotations $\Omega_{ij}$, the Jacobian of the map has spectral radius $1.173$, while the bound $\max_i\lambda_{\max}(M_i)/[\alpha+\lambda_{\max}(M_i)]$ with $M_i=\sum_j\beta_{ij}\Lambda_{t,ij}$ equals $0.9993$; iterating the map amplifies the error by a factor of about $10^3$ in $60$ steps, so it is not a contraction in any norm. Relaxing the restrictions removes the contraction proof, but it does not remove the canonical scalar potential. Strong monotonicity can instead provide a sufficient uniqueness criterion for the receiver-only field~\eqref{eq:receiver_only_field} on a bounded invariant set; that statement concerns the optional asymmetric approximation and not descent on the canonical scalar.

For the prospective protocol of Section~\ref{sec:participatory}, the consequence is that the frozen-model-channel arm ($\gamma_{ij}=0$) and the arm with declared nonzero slow coupling both retain the fast-channel belief alignment, so neither reduces to the exactly solvable cross-scale mean field. The open items are a perturbative bound on how the within-scale corrections move the Gaussian cross-scale fixed point, compatibility of cyclic row models, and uniqueness and convergence for the fully coupled, state-dependent $\mathrm{GL}^+(K_q)$ system.

\section{Beyond Mean Field: Correlated Posteriors, the Nested Bound, and Fluctuation Corrections}
\label{sec:beyond_mean_field_gauge_vfe}

Mean field enters the present construction in four places. The factorization $Q=\prod_i q_i$ suppresses state-posterior correlations. Deterministic minimization of one population functional selects a single configuration and suppresses configuration fluctuations. The exchangeable large-$N$ description replaces a finite labeled population by an empirical law. The operational coarse-graining construction projects an exact pushforward back into a selected multi-agent functional class. These are four different approximations, and improving one does not remove the other three.

The configuration-space Gibbs lift of Theorem~\ref{thm:belief_configuration_gibbs_lift} places a probability law on whole belief configurations and introduces the distinct configuration temperature $T_{\mathrm{cfg}}$, which is a bookkeeping scale of that ensemble and not a physical temperature. The Laplace term in Eq.~\eqref{eq:rg_laplace} also retains a local Gaussian fluctuation determinant when internal cluster variables are eliminated. Those results do not yet define a correlated posterior over the latent states of different agents, a generating functional for connected configuration correlations, or a self-consistent two-point theory. This appendix states the nested bound that organizes these objects, the correlated state families that preserve evidence semantics, and an outlook on fluctuation corrections, without reclassifying the peer-KL functional as an ordinary state-level evidence bound.

\subsection{A Normalized Nested Variational Model}

Let $X$ denote a slow generative configuration of prior and transition parameters, supports, graph or source priors, precision fields, and frame or edge variables. It does not contain the variational posterior densities $q_i$ or $s_i$. Let $Y=(Y_1,\ldots,Y_N)$ denote the state-level and model-level latent variables, and let $o$ denote the observations. The state model is a fixed normalized conditional law
\begin{equation}
p_\theta(dY,o\mid X)
:=p_\theta(dY\mid X)p_\theta(o\mid Y,X),
\label{eq:beyond_mf_state_joint}
\end{equation}
whose factors do not depend on the variational posterior used to approximate them. For a possibly correlated conditional posterior $Q_X(dY)$, define the state variational free energy in nats by
\begin{equation}
\mathcal{F}_{\mathrm{state}}[Q_X;X,o]
:=D_{\mathrm{KL}}\bigl(Q_X\|p_\theta(\cdot\mid X)\bigr)
-\mathbb{E}_{Q_X}\log p_\theta(o\mid Y,X).
\label{eq:beyond_mf_state_vfe}
\end{equation}
It obeys
\begin{equation}
\mathcal{F}_{\mathrm{state}}[Q_X;X,o]
=-\log p_\theta(o\mid X)
+D_{\mathrm{KL}}\bigl(Q_X\|p_\theta(\cdot\mid o,X)\bigr),
\label{eq:beyond_mf_state_elbo_identity}
\end{equation}
so correlated inference changes the approximation family without changing the evidence semantics.

\begin{theorem}[Nested state and configuration VFE]
\label{thm:nested_state_configuration_vfe}
Let $P_0(dX)$ be a proper configuration prior, let Eq.~\eqref{eq:beyond_mf_state_joint} be normalized for $P_0$-almost every $X$, and define $p_\theta(o):=\int p_\theta(o\mid X)P_0(dX)\in(0,\infty)$. Let $R\ll P_0$ be a configuration posterior and let $X\mapsto Q_X$ be a measurable probability kernel with $Q_X\ll p_\theta(\cdot\mid X)$ for $R$-almost every $X$. Assume the terms below are finite, or interpret the identity in the extended-real sense without an indeterminate $\infty-\infty$. Then
\begin{align}
\mathcal{J}[R,Q]
&:=T_{\mathrm{cfg}}D_{\mathrm{KL}}(R\|P_0)
+T_{\mathrm{cfg}}\int
\mathcal{F}_{\mathrm{state}}[Q_X;X,o]R(dX)
\nonumber\\
&=-T_{\mathrm{cfg}}\log p_\theta(o)
+T_{\mathrm{cfg}}D_{\mathrm{KL}}
\bigl(R(dX)Q_X(dY)\|P_\theta(dX,dY\mid o)\bigr).
\label{eq:nested_state_configuration_identity}
\end{align}
Consequently, $-\mathcal{J}/T_{\mathrm{cfg}}$ is an evidence lower bound for one normalized hierarchical model, with equality if and only if the joint variational law equals the exact posterior almost everywhere.
\end{theorem}

\begin{proof}
Expand the joint KL divergence by the chain rule, use $P_\theta(dX,dY,o)=P_0(dX)p_\theta(dY,o\mid X)$, and collect the conditional terms into Eq.~\eqref{eq:beyond_mf_state_vfe}.
\end{proof}

The quantity $\mathcal{J}/T_{\mathrm{cfg}}$ is the two-level case, with latent $(X,Y)$, of the tower variational free energy of the companion theory~\cite[Ch.~12]{DennisGaugeVFE2026}. It is stated here because the configuration temperature and the vacuum prior below are specific to this paper. If $P_0$ is itself defined by a vacuum functional,
\begin{equation}
P_0(dX)
:=Z_0^{-1}\exp[-\mathcal{F}_{\mathrm{vac}}(X)/T_{\mathrm{cfg}}]\nu_0(dX),
\label{eq:beyond_mf_configuration_prior}
\end{equation}
then the first term of Eq.~\eqref{eq:nested_state_configuration_identity} is, up to the constant $T_{\mathrm{cfg}}\log Z_0$, the sum of $\mathbb{E}_R\mathcal{F}_{\mathrm{vac}}$ and $T_{\mathrm{cfg}}D_{\mathrm{KL}}(R\|\nu_0)$. This nests state-level evidence inside the configuration ensemble. The vacuum functional may contain only observation-free structural energies in the declared variables $X$; it cannot contain live $q_i$, $s_i$, or their peer-KL terms without promoting separate measure-valued fields to generative variables and specifying their relation to $Q_X$. The held-out likelihood is excluded from $\mathcal{F}_{\mathrm{vac}}$ and enters once through Eq.~\eqref{eq:beyond_mf_state_vfe}. Environmental variables may be represented as ordinary latent nodes or observed boundary conditions; the construction does not treat a variational belief of another agent as a fixed generative factor.

Both terms of $\mathcal{J}$ carry $T_{\mathrm{cfg}}$, so at the optimal state posterior the optimal configuration posterior is $R^*(dX)\propto\exp[-\mathcal{F}_{\mathrm{vac}}(X)/T_{\mathrm{cfg}}]p_\theta(o\mid X)\nu_0(dX)$. As $T_{\mathrm{cfg}}\to0$ it concentrates on the minimizers of $\mathcal{F}_{\mathrm{vac}}$ independently of $o$: for $\mathcal{F}_{\mathrm{vac}}(X)=X^2/2$ and $p_\theta(o\mid X)=\mathcal{N}(o;X,1)$ one finds $\mathbb{E}_{R^*}[X]=oT_{\mathrm{cfg}}/(1+T_{\mathrm{cfg}})$. With this scaling the data enter the configuration law only at order $T_{\mathrm{cfg}}$, the same order as the fluctuation corrections below, so the limit $T_{\mathrm{cfg}}\to0$ does not recover data-driven minimization of the multi-agent functional.

\subsection{Correlated State Posteriors and Region Free Energies}

The first extension beyond mean field preserves the normalized model above while enlarging $Q_X$. Suppose the conditional state prior has a gauge-covariant pairwise factorization on a graph $\mathcal{G}=(V,E)$,
\begin{equation}
p_\theta(Y\mid X)
:=Z_X^{-1}
\prod_{i\in V}\psi_i(Y_i;X)
\prod_{(i,j)\in E}\psi_{ij}(Y_i,Y_j;X),
\label{eq:beyond_mf_pairwise_joint}
\end{equation}
where each edge factor is a fixed model potential and not a function of live variational densities. Introduce node beliefs $b_i$ and edge beliefs $b_{ij}$ satisfying normalization and local marginal consistency. Their Bethe entropy is
\begin{equation}
H_{\mathrm{B}}[b]
:=\sum_{(i,j)\in E}H[b_{ij}]
-\sum_{i\in V}(d_i-1)H[b_i].
\label{eq:beyond_mf_bethe_entropy}
\end{equation}
The corresponding region free energy is the expectation of the fixed node and edge energies $-\log\psi_i$ and $-\log\psi_{ij}$ minus $H_{\mathrm{B}}$. On a tree the local marginal polytope equals the marginal polytope~\cite[Prop.~4.1]{Wainwright2008}, and minimizing this functional over it returns the exact marginals and $-\log Z_X$. On a loopy graph it is the Bethe approximation. Larger regions give the cluster-variation method~\cite{Kikuchi1951} and its region-graph generalizations, a family of more expressive approximations whose accuracy depends on the choice of regions and is not guaranteed to improve~\cite[Sec.~VIII.C]{Yedidia2005}. Generic loopy pseudomarginals need not be the marginals of any global $Q_X$, and their region objective is not a generic upper or lower evidence bound. Bound-preserving constructions require additional structure, such as a globally realizable variational family or an appropriate tree-reweighted relaxation. Ordinary mean field is recovered by the restriction $b_{ij}=b_i b_j$ on every edge when those beliefs arise from a global product law. It is therefore a lower-complexity variational subfamily of the correlated theory, not the definition of the theory.

Gauge covariance acts on the region beliefs as a joint pushforward. If $Y_i\mapsto G_iY_i$ and $Y_j\mapsto G_jY_j$, then $b_{ij}$ is pushed forward by $G_i\times G_j$, while the edge potential is transformed contragrediently so that its scalar energy is unchanged. Entropies must be defined relative to a covariantly transformed or quasi-invariant reference measure, including its Jacobian; a finite invariant reference is not generally available for noncompact $\mathrm{GL}^+(K_q)$. A Regime~I link $\Omega_{ij}=U_iU_j^{-1}$ has trivial loop holonomy and is removed by a passive frame re-expression. A correlated theory with nontrivial holonomy must instead retain the independent links on the interaction graph of Regime~II (Section~\ref{sec:discrete_regime_ii}), with orientation reversal imposed; integrating only vertex frames cannot generate nontrivial holonomy.

For the Gaussian sector, a computationally direct correlated family is
\begin{equation}
Q_X(Y)=\mathcal{N}(\bar{Y},C),
\qquad
C^{-1}=\Lambda_Q,
\label{eq:beyond_mf_block_gaussian}
\end{equation}
with a sparse block precision $\Lambda_Q$ supported on the state interaction graph. The off-diagonal blocks $C_{ij}$ carry inter-agent correlations, and if $Y_i=(k_i,m_i)$ contains fast state and slow model variables, the within-node block $C_i^{(km)}$ carries cross-channel correlations that the separate $q_i$ and $s_i$ factors discard. Under a local linear frame change, $C_{ij}\mapsto \rho_i(G_i)C_{ij}\rho_j(G_j)^\top$, where $\rho_i$ is block diagonal if the state and model channels use distinct representations, so the covariance family is closed under the group action. The block-diagonal restriction $C_{ij}=0$ for $i\ne j$ is the Gaussian mean-field limit. At a Gaussian target of precision $\Lambda$, its exact cost at the optimum is the determinant gap $\frac12[\sum_b\log\det\Lambda_{bb}-\log\det\Lambda]\ge0$, which vanishes exactly when all cross-block precisions vanish~\cite[Ch.~G4]{DennisGaugeVFESuppG2026}. Sparse block-Gaussian inference is therefore the smallest extension that exposes a measurable connected two-point state function while retaining tractable linear algebra.

The live alignment term $D_{\mathrm{KL}}(q_i\|\Omega_{ij}q_j)$ does not become an exact instance of Eq.~\eqref{eq:beyond_mf_pairwise_joint} merely because edge beliefs have been introduced: its target still moves with another variational factor, which is the obstruction proved by Theorem~\ref{thm:state_level_elbo_nogo}. It has two coherent placements. It may remain a configuration-level consensus potential, or it may be replaced by a fixed state-transition or pair potential from which the correlated posterior is inferred. Using both for the same interaction would count that edge twice.

\begin{remark}[No-double-counting conditions]
\label{rem:beyond_mf_no_double_counting}
An observation likelihood enters the state free energy once and is not also placed in the vacuum action; a physical interaction is represented either by a fixed edge factor of the state model or by a consensus potential on configurations, not both; attention rows are either integrated as random simplex variables or eliminated at their envelope optimum $\beta^*$, and the row entropy is not added again after elimination; and marginal Gaussian entropies and the correlation correction are components of one state-posterior entropy, while attention-row entropy and configuration meta-entropy belong to separate sample spaces. These are instances of two general statements of the companion theory: an energy appended to a joint without its normalizer is not an evidence bound~\cite[Ch.~4]{DennisGaugeVFE2026}, and the complexity of a correlated recognition law splits exactly into its total correlation and its marginal divergences~\cite[Ch.~5]{DennisGaugeVFE2026}.
\end{remark}

\subsection{Outlook: Fluctuation Corrections and the Effective Action}

After the state family has been specified, define a state-reduced configuration action
\begin{equation}
\mathcal{A}_o[X]
:=\mathcal{F}_{\mathrm{vac}}[X]
+T_{\mathrm{cfg}}
\inf_{Q_X\in\mathcal{Q}_{\mathrm{corr}}}
\mathcal{F}_{\mathrm{state}}[Q_X;X,o].
\label{eq:beyond_mf_bare_action}
\end{equation}
With the unrestricted posterior family the second term is $-T_{\mathrm{cfg}}\log p_\theta(o\mid X)$; with a normalized family such as a sparse block-Gaussian $Q_X$ it is a variational upper approximation to that conditional surprisal, while a Bethe or Kikuchi pseudomarginal optimization does not inherit this evidence-bound statement. The live multi-agent functional cannot serve as $\mathcal{F}_{\mathrm{vac}}$ in Theorem~\ref{thm:nested_state_configuration_vfe}, because it depends on the variational $q_i,s_i$ and contains the observation term. The peer-KL sector is a gauge-covariant consensus energy and, when its belief-configuration partition function is finite, an exact generalized-Bayesian energy on that separate sample space; inside $\mathcal{A}_o$ it is only a candidate truncation pending an explicit summary map from $Q_X$ to belief fields and quantitative energy and force residuals.

Let $\mathfrak{X}_{\mathrm{phys}}$ be the configuration space after quotienting only the transformations declared to be redundancies; active frame changes remain coordinates of $X$. Around a background $\bar{X}$, write $X=\operatorname{Exp}_{\bar{X}}\xi$ in normal coordinates on the quotient and define
\begin{equation}
Z_{\bar{X}}[J]
:=\int_{\mathfrak{X}_{\mathrm{phys}}}
\exp\left[-\frac{\mathcal{A}_o[\operatorname{Exp}_{\bar{X}}\xi]-J_A\xi^A}{T_{\mathrm{cfg}}}\right]
d\nu_{\mathrm{phys}}(\xi),
\qquad
W_{\bar{X}}[J]:=T_{\mathrm{cfg}}\log Z_{\bar{X}}[J].
\label{eq:beyond_mf_generating_functional}
\end{equation}
The mean displacement is $\varphi^A=\delta W/\delta J_A$, and the Legendre transform $\Gamma_{\bar{X}}[\varphi]:=J_A\varphi^A-W_{\bar{X}}[J]$ is the one-particle-irreducible (1PI) effective action in the chosen background chart. With this sign convention its value at the source-free mean is $-T_{\mathrm{cfg}}\log Z_{\bar{X}}[0]$, the configuration free energy, and its Hessian gives the connected configuration covariance $G^{AB}=T_{\mathrm{cfg}}[(\Gamma^{(2)})^{-1}]^{AB}$ on the physical modes, a covariance across whole configurations distinct from the state covariance $C_{ij}$ of Eq.~\eqref{eq:beyond_mf_block_gaussian}. The 1PI action is not the coarse energy $\mathcal{F}^{\mathrm{exact}}_{\ell+1}$ of Appendix~\ref{sec:meta_agent_rg}, which is a conditional partition function of eliminated variables (a Wilsonian effective energy in the block-variable sense, without rescaling), and it is background- and base-point-dependent until compatible constructions on overlapping quotient charts are supplied. Write $d\nu_{\mathrm{phys}}=\rho_{\mathrm{phys}}\operatorname{vol}_{g_{\mathrm{phys}}}$ and assume that the ambient metric and reference measure are invariant under the quotiented redundancy. When $\bar{X}$ is a nondegenerate saddle on the quotient and the covariant Hessian $H_{\mathrm{phys}}$ of $\mathcal{A}_o$ on the physical tangent modes is positive, the Laplace expansion gives
\begin{equation}
\Gamma[\bar{X}]
=\mathcal{A}_o[\bar{X}]
+\frac{T_{\mathrm{cfg}}}{2}
\operatorname{Tr}\log\left(\frac{g_{\mathrm{phys}}^{-1}H_{\mathrm{phys}}}{2\pi T_{\mathrm{cfg}}}\right)
+\Gamma_{\mathrm{measure}}[\bar{X}]
+O(T_{\mathrm{cfg}}^2),
\qquad
\Gamma_{\mathrm{measure}}=-T_{\mathrm{cfg}}\log\rho_{\mathrm{phys}}(\bar{X}).
\label{eq:beyond_mf_one_loop}
\end{equation}
The factor $2\pi T_{\mathrm{cfg}}$ is part of the result; replacing it by a fixed scale leaves an $O(T_{\mathrm{cfg}}\log T_{\mathrm{cfg}})$ remainder. For a free action of the redundancy group by isometries of the ambient metric $g$, with $\nu$ the Riemannian volume of $g$, that volume factors locally into the quotient volume and the orbit volume, so $\rho_{\mathrm{phys}}=(\det\mathcal{G}_{ab})^{1/2}$ with $\mathcal{G}_{ab}=g(V_a,V_b)$ the Gram matrix of the redundancy generators, and $\Gamma_{\mathrm{measure}}=-\frac{T_{\mathrm{cfg}}}{2}\log\det\mathcal{G}$. The Frobenius frame metric of Eq.~\eqref{eq:pullback_metric} is not invariant under the right redundancy $U_i\mapsto U_ig$, so for $\mathrm{GL}^+(K_q)$ frames $g_{\mathrm{phys}}$ is defined only after a right-invariant metric is adopted. An equivalent redundant-field calculation instead imposes a gauge condition $\Phi_{\mathrm{gf}}=0$ and includes one Faddeev--Popov factor~\cite{FaddeevPopov1967}: in coordinates orthonormal for $g$ at $\bar{X}$, the one-loop and measure terms together equal $\frac{T_{\mathrm{cfg}}}{2}\log\det(H_{\mathrm{slice}}/2\pi T_{\mathrm{cfg}})-T_{\mathrm{cfg}}\log|\det M_{\mathrm{FP}}|$, where $H_{\mathrm{slice}}$ is the Hessian on the slice and $M_{\mathrm{FP}}=\partial\Phi_{\mathrm{gf}}/\partial\theta$, because $\det M_{\mathrm{FP}}^2/\det H_{\mathrm{slice}}=\det\mathcal{G}/\det H_{\mathrm{phys}}$. The direct-quotient and gauge-fixed prescriptions are alternatives and must not be multiplied together. A background-field construction is needed if the redundant-field expansion is to retain covariance under background gauge transformations~\cite{DeWitt1967,Abbott1981,Abbott1982}. Raw Haar integration over noncompact $\mathrm{GL}^+(K_q)$ is not admissible: redundant noncompact directions require a quotient or gauge fixing, and noncompact active directions require a proper confining reference measure, so that Eq.~\eqref{eq:beyond_mf_generating_functional} exists. Equation~\eqref{eq:beyond_mf_one_loop} extends the cluster-specific log determinant in Eq.~\eqref{eq:rg_laplace} to the full physical configuration Hessian; it is a local asymptotic expansion, not a nonperturbative definition when the Hessian has additional soft or unstable modes.

The 1PI action closes equations for the mean configuration but treats $G$ as the inverse Hessian at the selected truncation. A two-particle-irreducible (2PI) extension would add a symmetric bilocal source $K_{AB}$ and perform a double Legendre transform to $\Gamma_{\mathrm{2PI}}[\varphi,G]$, whose stationarity condition $\delta\Gamma_{\mathrm{2PI}}/\delta G=0$ gives a self-consistent configuration covariance~\cite{CornwallJackiwTomboulis1974}. The gauge-consistent 2PI construction is a proposed completion, not a result derived from the five-term functional; it would matter when soft collective modes make the one-loop determinant large or when correlation feedback moves the saddle or changes its stability.

As Appendix~\ref{sec:meta_agent_rg} records, the exact pushforward generally generates every operator allowed by the retained symmetries, here including higher-order region interactions, nonlocal kernels, covariance couplings, curvature terms, and measure contributions~\cite{Wilson1974}; closure in the original five-term family is therefore a truncation hypothesis, to be tested by the calibrated energy residual $\delta_{I,0}$ and the force residual $\delta_{I,1}$ of Eq.~\eqref{eq:rg_closure}, not assumed from gauge covariance. A complete calculation would compare at least three nested truncations: the present mean-field saddle, a correlated state family with one-loop configuration corrections, and a 2PI or region-enlarged theory in which the two-point functions are varied self-consistently; the change in held-out evidence, connected correlations, spectra, and closure residuals decides whether the extra operators earn their place. This ordering defines a research program, not a controlled approximation ladder. Mean field is formally recovered by factorizing the state posterior. Exchangeable empirical-measure closure and projection of the coarse-grained image back to the five-term family are formal limits or truncations until propagation-of-chaos and quantitative closure-residual hypotheses are established.
